% !TEX encoding = UTF-8
% !TEX TS-program = pdflatex
% !TEX root = ../tesi.tex

%**************************************************************
\chapter{Introduzione}
\label{cap:introduzione}
%**************************************************************

% Introduzione al contesto applicativo.\\

% \noindent Esempio di utilizzo di un termine nel glossario \\
% \gls{api}. \\

% \noindent Esempio di citazione in linea \\
% \cite{site:agile-manifesto}. \\

% \noindent Esempio di citazione nel pie' di pagina \\
%citazione\footcite{womak:lean-thinking} \\

%**************************************************************
\section{L'azienda}

\href{https://scaiitec.it/}{SCAI ITEC} è un'azienda italiana appartenente al \href{https://www.grupposcai.it/}{gruppo SCAI}. ITEC si occupa di consulenza, \textit{System Integration} e \textit{Application management} in ambito ICT. L'azienda opera principalmente nel settore bancario, assicurativo, industriale e di pubblica amministrazione e servizi. \\Gli elementi chiave del successo e della crescita di SCAI ITEC sono:
\begin{itemize}
    \item vasta e profonda conoscenza delle tecnologie;
    \item grande attenzione per la soddisfazione del cliente;
    \item molta esperienza, maturata nel corso del tempo.
\end{itemize} 
ITEC in ambito ICT è oggi una delle maggiori realtà del nord-est del Paese, proponendosi come partner per qualsiasi tipo di applicazione, di progetto e di servizio modellato sulle specifiche necessità del cliente. 
\\Grazie alla consolidata esperienza nel ruolo di \textit{\gls{systemIntegratorg}\glsfirstoccur{}} e alle soluzioni \textit{leader} di mercato proposte, ITEC è in grado di garantire ai propri clienti risposte rapide, concrete e qualificate in base alle specifiche esigenze di tipo gestionale e applicativo.
\\Ultimo, ma non meno importante, dei motivi per cui l'azienda è all'avanguardia rispetto alle nuove tecnologie è il grande investimento di ITEC in ricerca, sviluppo e formazione del personale.

\begin{figure}[h]
    \begin{center}
    \includegraphics[width=0.5\textwidth]{logo-itec}
    \caption{Logo SCAI ITEC}
    \label{fig:figure1}
    \end{center}
\end{figure}


%**************************************************************
\section{Scopo dello stage}
\label{sez:scopo}
La finalità principale dello stage è stato quella di inserire lo studente all'interno di una nuova progettualità, con un particolare \textit{focus} sulle tematiche legate alle tecnologie multimediali e alla loro distribuzione. Lo studente, affiancato da un \textit{IT Architect}, ha avuto la possibilità di partecipare al disegno, alla progettazione e alla realizzazione di una nuova progettualità. L'obiettivo è far apprendere le tecnologie e le \textit{best practice} utilizzate in azienda durante il ciclo di vita di un applicativo.
La progettualità vista è volta a creare un \textit{software} per la gestione di \gls{contenutog}\glsfirstoccur{}, per la loro creazione, modifica e distribuzione nei vari canali di vendita.
Grazie a questa nuova applicazione si potrà:
\begin{itemize}
    \item creare, modificare, eliminare e clonare i \gls{contenutog};
    \item raggruppare i \gls{contenutog} su segmenti di mercato e distribuirli;
    \item pianificare la visualizzazione e l’aggiornamento dei \gls{contenutog} sui vari canali di distribuzione;
    \item seguire un processo di \textit{Authoring} (paradigma Editore, Redattore, Supervisore) nella fase di creazione e distribuzione.
\end{itemize} 
In quanto l'intero progetto è di dimensione non indifferente, l'obiettivo dello stage è stato di sviluppare le funzionalità relative al ruolo di editore, più nel dettaglio:
\begin{itemize}
    \item creazione di un contenuto;
    \item modifica di un contenuto;
    \item eliminazione di un contenuto;
    \item preparazione di un contenuto per la distribuzione;
    \item \textit{auditing} delle azioni effettuate dagli utenti;
    \item documentazione delle funzionalità implementate.
\end{itemize}

\section{Criticità e soluzioni}
Durante la pianificazione sono state evidenziate alcune criticità. Queste vengono descritte di seguito insieme alle soluzioni trovate per mitigarle.

\subsection*{Conoscenza delle tecnologie}
Il progetto di stage ha previsto l'utilizzo di molteplici tecnologie mai utilizzate da me fino a quel momento. Questo ha introdotto il rischio di un eccessivo \textit{overhead} per l'apprendimento delle stesse. La soluzione adottata è stata la seguente:
\begin{itemize}
    \item ogni tecnologia mai utilizzata durante il mio percorso di studi è stata introdotta brevemente dal tutor aziendale;
    \item mi è stato fornito materiale da consultare per apprendere le basi delle nuove tecnologie;
    \item è stato infine deciso di utilizzare l'approccio \textit{"Learn by doing"} con la possibilità di chiedere aiuto al tutor aziendale nel caso fosse necessario.
\end{itemize} 

\subsection*{Coordinamento del team}
A causa del COVID-19 è stato deciso di comune accordo di svolgere il lavoro in modalità telematica per tutta la durata dello stage. Questo ha introdotto il rischio che ci fossero maggiori difficoltà nelle interazioni con il tutor aziendale e con il team di sviluppo. Tale eventualità è stata mitigata grazie all'applicazione di \hyperref[sez:scrum]{Scrum}, che impone vengano effettuati dei meeting giornalieri per favorire il coordinamento del team.

\subsection*{Progettazione del database}
Uno dei momenti più importanti della fase di \hyperref[cap:analisi]{Analisi e Progettazione} è stata la progettazione della base di dati. Un'ideazione errata del \textit{database} avrebbe infatti introdotto numerosi problemi durante la codifica, con il rischio di dover ripartire da zero nello sviluppo. Per mitigare tale rischio la base di dati una volta progettata è stata controllata attentamente da parte del tutor aziendale.

\section{Organizzazione del testo}

\begin{description}
    \item[{\hyperref[cap:obiettivi-pianificazione]{Il secondo capitolo}}] descrive gli obiettivi dello stage, la pianificazione del lavoro effettuata a monte e le aspettative personali riguardanti lo stage;
    
    \item[{\hyperref[cap:tecnologie]{Il terzo capitolo}}] descrive le tecnologie e gli strumenti di sviluppo utilizzati durante il corso dello stage;
    
    \item[{\hyperref[cap:metodologia-lavoro]{Il quarto capitolo}}] approfondisce la metodologia di lavoro e i ruoli adottati dal team di sviluppo;
    
    \item[{\hyperref[cap:analisi]{Il quinto capitolo}}] descrive l'analisi, la progettazione iniziale e le soluzioni adottate durante la codifica;
    
    \item[{\hyperref[cap:prodotto]{Il sesto capitolo}}] descrive le funzionalità del prodotto finale, la documentazione e i test eseguiti;
    
    \item[{\hyperref[cap:considerazioni]{Il settimo capitolo}}] contiene le considerazioni finali riguardanti lo stage e una valutazione personale sul lavoro svolto;

\end{description}

Riguardo la stesura del testo, relativamente al documento, sono state adottate le seguenti convenzioni tipografiche:
\begin{itemize}
	\item gli acronimi, le abbreviazioni e i termini ambigui o di uso non comune menzionati vengono definiti nel glossario, situato alla fine del presente documento;
	\item per la prima occorrenza dei termini riportati nel glossario viene utilizzata la seguente nomenclatura: \emph{parola}\glsfirstoccur{};
	\item i termini in lingua straniera o facenti parti del gergo tecnico sono evidenziati con il carattere \emph{corsivo}.
\end{itemize}